\section{Conclusion}
\label{sec:conclusion}

This project implemented ring-allreduce from scratch directly on MPI point-to-point primitives, with the two-phase reduce-scatter and all-gather structure, its edge cases, and its bandwidth-optimal data volume all treated as first-class requirements rather than details to fill in later, and validated against both vendor \texttt{MPI\_Allreduce} and an independently-computed reference across a test matrix that deliberately includes non-power-of-two process counts and the zero-size-chunk cases a smaller test matrix would miss. The analysis pipeline fits this project's own hardware constants, $\alpha$ and $\beta$, from measured data via a weighted least-squares procedure whose necessity (and whose alternative's failure mode) was itself demonstrated empirically rather than assumed, and uses those constants to separate pipeline fill latency from rank synchronization overhead as two distinct, individually quantified small-message bottlenecks, and to explain, via step count rather than vague appeal to "vendor optimization," the specific circumstances under which this implementation can lose to a vendor \texttt{MPI\_Allreduce}'s own algorithm selection despite being bandwidth-optimal. On the single-node Microsoft MPI measurements analyzed here, that loss is real below a few kilobytes and reverses above it, the ring becoming the faster of the two by up to a factor of two through the cache-resident message sizes, which is the expected consequence of a bandwidth-optimal but step-heavy algorithm rather than a defect.

The limitations, stated plainly rather than buried, are those of a single-node run. The transport underneath is Microsoft MPI's shared memory, not a network fabric, which is what produces the cache-resident bus-bandwidth peak and the modest fit of the linear cost model (Section~\ref{sec:dataset-note}); and Microsoft MPI's lack of \texttt{coll/tuned} controls means the forced ring-versus-ring comparison, which would isolate implementation quality from algorithm choice, was not among the variants that could be measured. Both are addressed by the same step, an Open~MPI multi-node sweep via \texttt{scripts/run\_slurm\_sweep.sbatch}, regenerating this report from that output (\texttt{make report} against the new CSVs) without changing a single line of the methodology described here.

Beyond that, three extensions would most directly build on this project's own findings rather than add unrelated scope. First, and most directly motivated by Section~\ref{sec:discussion}'s own data: a segmented, pipelined ring that overlaps communication across steps, attacking the latency term this report's decomposition specifically identified as the one still worth optimizing. Second, a naive gather-to-one-rank-then-broadcast baseline, purely as a third comparison point that would make concrete, in this project's own numbers, exactly why a from-scratch ring was worth building in the first place. Third, a real multi-node sweep via \texttt{scripts/run\_slurm\_sweep.sbatch}, extending every method in this report to a setting where network fabric, rather than shared memory, sets the real $\alpha$ and $\beta$. A GPU and NCCL comparison is deliberately not on this list: it is a different hardware regime with its own transfer costs and topology effects that a CPU MPI benchmark cannot speak to responsibly, and is noted here only as future work rather than attempted with a mismatched proxy.
